\documentclass[11pt,a4paper]{article}
\usepackage[utf8]{inputenc}
\usepackage{amsmath,amssymb,amsthm}
\usepackage{algorithm}
\usepackage{algpseudocode}
\usepackage{booktabs}
\usepackage{hyperref}
\usepackage{geometry}
\geometry{margin=1in}

\title{Receptive Field Analysis: Theoretical Recurrence, Coordinate Mapping, and Gaussian Effective Distribution}
\author{Team Scaibu \\ \small Email: \texttt{jaydeep.9664920749@gmail.com} \\ \small Architecture and Core Engine Team}
\date{\today}

\newtheorem{theorem}{Theorem}
\newtheorem{lemma}[theorem]{Lemma}
\newtheorem{proposition}[theorem]{Proposition}
\theoremstyle{definition}
\newtheorem{definition}{Definition}
\newtheorem{remark}{Remark}

\begin{document}

\maketitle

\begin{abstract}
The receptive field of a unit in a deep convolutional neural network defines the spatial extent of input pixels that can mathematically influence its activation. While the theoretical receptive field (TRF) grows linearly with layer depth and dilation factors, the effective receptive field (ERF)—defined by the gradient distribution of the output with respect to input pixels—decays rapidly away from the center, following a two-dimensional Gaussian profile according to the Central Limit Theorem. This document formalizes the discrete linear recurrences for TRF size, cumulative stride (jump), spatial coordinate anchoring, and derives the closed-form variance of the Gaussian ERF.
\end{abstract}

\section{Notation and Mathematical Preliminaries}

Let a convolutional network be parameterized by an ordered sequence of $L$ layers $\mathcal{L}_1, \dots, \mathcal{L}_L$, where each layer $\mathcal{L}_l$ is defined by the 4-tuple $(k_l, s_l, p_l, d_l)$ representing kernel size, stride, padding, and dilation factor respectively.

\begin{table}[h]
\centering
\caption{Mathematical Notation for Receptive Field Analysis}
\begin{tabular}{lll}
\toprule
\textbf{Symbol} & \textbf{Domain} & \textbf{Semantic Description} \\
\midrule
$k_l'$ & $\mathbb{N}_{\ge 1}$ & Dilated effective kernel size: $k_l' = d_l (k_l - 1) + 1$ \\
$s_l$ & $\mathbb{N}_{\ge 1}$ & Stride of layer $l$ \\
$p_l$ & $\mathbb{N}_{\ge 0}$ & Padding of layer $l$ \\
$r_l$ & $\mathbb{R}_{\ge 1}$ & Theoretical Receptive Field size at layer $l$ ($r_0 = 1$) \\
$j_l$ & $\mathbb{R}_{\ge 1}$ & Cumulative stride (jump) at layer $l$ ($j_0 = 1$) \\
$\text{start}_l$ & $\mathbb{R}$ & Center position of index 0 in input coordinate space ($\text{start}_0 = 0.5$) \\
$\sigma_{\text{ERF}}^2$ & $\mathbb{R}_{> 0}$ & Spatial variance of the Gaussian Effective Receptive Field \\
\bottomrule
\end{tabular}
\end{table}

\section{Problem Definition and Core Concepts}

\subsection{3-Level Explanation}
\begin{enumerate}
    \item \textbf{Intuitive (High School level):} If you hold a magnifying glass, looking through multiple lenses makes each spot in the final image see a wider patch of the original object. Receptive field analysis tells us exactly how big that patch is and where its center points on the original picture.
    \item \textbf{Engineering Level:} Object detectors and segmentors must match their anchor box dimensions and feature levels to the physical receptive field of network backbones. If a small feature map unit has a theoretical receptive field of $300 \times 300$ pixels, an anchor of size $32 \times 32$ will suffer from contextual mismatch.
    \item \textbf{Mathematical Level:} The theoretical receptive field $r_L$ is the diameter of the support $\text{supp}\left( \frac{\partial y_L}{\partial x_0} \right)$. By linear expansion through convolutional convolution operators, the support grows according to the discrete recurrence $r_l = r_{l-1} + (k_l' - 1) j_{l-1}$.
\end{enumerate}

\subsection{5-Step Algorithmic Framework}
\begin{enumerate}
    \item \textbf{Kernel Dilation Adjustment:} Compute effective kernel diameter $k_l' = d_l(k_l - 1) + 1$.
    \item \textbf{Theoretical Size Recurrence:} Update receptive field size $r_l = r_{l-1} + (k_l' - 1) j_{l-1}$.
    \item \textbf{Coordinate Anchor Offset:} Update center of index 0: $\text{start}_l = \text{start}_{l-1} + \left(\frac{k_l' - 1}{2} - p_l\right) j_{l-1}$.
    \item \textbf{Cumulative Stride Compounding:} Update cumulative stride: $j_l = j_{l-1} \cdot s_l$.
    \item \textbf{Effective Gaussian Variance Calculation:} Accumulate per-layer kernel variances to compute the Gaussian ERF envelope.
\end{enumerate}

\section{Algorithm Specification}

\begin{algorithm}[H]
\caption{Receptive Field and Spatial Coordinate Profiling}
\label{alg:rf_analysis}
\begin{algorithmic}[1]
\Require Layer specifications $\{(k_l, s_l, p_l, d_l)\}_{l=1}^L$.
\Ensure Theoretical RF sizes $\{r_l\}$, jumps $\{j_l\}$, start offsets $\{\text{start}_l\}$, ERF standard deviation $\sigma_{\text{ERF}}$.
\State $r_0 \gets 1.0, \quad j_0 \gets 1.0, \quad \text{start}_0 \gets 0.5, \quad \text{var}_{\text{sum}} \gets 0.0$
\For{$l = 1$ \textbf{to} $L$}
    \State $k_l' \gets d_l (k_l - 1) + 1$
    \State $r_l \gets r_{l-1} + (k_l' - 1) \cdot j_{l-1}$
    \State $\text{start}_l \gets \text{start}_{l-1} + \left( \frac{k_l' - 1}{2} - p_l \right) \cdot j_{l-1}$
    \State $\text{var}_{\text{sum}} \gets \text{var}_{\text{sum}} + \left( \frac{k_l^2 - 1}{12} \right) \cdot j_{l-1}^2$
    \State $j_l \gets j_{l-1} \cdot s_l$
\EndFor
\State $\sigma_{\text{ERF}} \gets \sqrt{\text{var}_{\text{sum}}}$
\State \Return $\{r_l\}_{l=1}^L, \{j_l\}_{l=1}^L, \{\text{start}_l\}_{l=1}^L, \sigma_{\text{ERF}}$
\end{algorithmic}
\end{algorithm}

\section{Theoretical Guarantees and Central Limit Convergence}

\begin{theorem}[Gaussian Effective Receptive Field Theorem]
\label{thm:erf_gaussian}
Let $y$ denote the output activation at the center of an $L$-layer convolutional network with random i.i.d. zero-mean weights. Under the Central Limit Theorem for convolutions of independent spatial density functions, the normalized gradient distribution:
\begin{equation}
\mathcal{G}(u, v) = \frac{\left| \frac{\partial y}{\partial x(u, v)} \right|}{\sum_{u', v'} \left| \frac{\partial y}{\partial x(u', v')} \right|}
\end{equation}
converges weakly as $L \to \infty$ to an isotropic 2D Gaussian distribution $\mathcal{N}(0, \sigma^2 I)$ with variance:
\begin{equation}
\sigma^2 = \sum_{l=1}^L \frac{k_l^2 - 1}{12} \cdot j_{l-1}^2
\end{equation}
Consequently, the effective receptive field only covers a fraction $\mathcal{O}(1/\sqrt{L})$ of the theoretical receptive field $r_L$.
\end{theorem}

\begin{proof}
Each convolutional layer acts as a spatial discrete filtering operator with kernel variance $\sigma_l^2 = \frac{k_l^2 - 1}{12}$. In the continuous limit, the backward gradient path is equivalent to the convolution of $L$ uniform rectangular probability distributions. By the Classical Central Limit Theorem for independent random variables, the $L$-fold convolution of bounded continuous distributions converges uniformly to a Gaussian distribution whose total variance is the sum of the individual layer variances scaled by the spatial stride factors $j_{l-1}^2$.
\end{proof}

\section{Complexity Analysis}

\begin{itemize}
    \item \textbf{Time Complexity:} $\mathcal{O}(L)$ where $L$ is the number of layers in the CNN architecture (constant arithmetic operations per layer).
    \item \textbf{Space Complexity:} $\mathcal{O}(L)$ to store layer-by-layer profile tuples.
\end{itemize}

\section{Exactness and Error Analysis}

The discrete recurrence $r_l = r_{l-1} + (k_l' - 1)j_{l-1}$ is exact and closed-form. The coordinate mapping $\text{center}(i) = \text{start}_l + i \cdot j_l$ exactly identifies the continuous sub-pixel centroid of the receptive field in input pixel coordinates.

\section{Worked Numerical Example}

Consider a 2-layer CNN:
\begin{itemize}
    \item Layer 1: $k_1 = 3, s_1 = 2, p_1 = 1, d_1 = 1$.
    \item Layer 2: $k_2 = 3, s_2 = 1, p_2 = 1, d_2 = 1$.
\end{itemize}

\textbf{Layer 1 Calculation ($r_0 = 1, j_0 = 1, \text{start}_0 = 0.5$):}
\begin{itemize}
    \item $k_1' = 1(3-1)+1 = 3$.
    \item $r_1 = 1 + (3 - 1) \cdot 1 = 3$.
    \item $\text{start}_1 = 0.5 + \left( \frac{3-1}{2} - 1 \right) \cdot 1 = 0.5 + (1 - 1) = 0.5$.
    \item $j_1 = 1 \cdot 2 = 2$.
    \item $\text{var}_1 = \frac{3^2 - 1}{12} \cdot 1^2 = \frac{8}{12} = 0.6667$.
\end{itemize}

\textbf{Layer 2 Calculation:}
\begin{itemize}
    \item $k_2' = 3$.
    \item $r_2 = r_1 + (3 - 1) \cdot j_1 = 3 + 2 \cdot 2 = 7$.
    \item $\text{start}_2 = \text{start}_1 + \left( \frac{3-1}{2} - 1 \right) \cdot j_1 = 0.5 + (0) \cdot 2 = 0.5$.
    \item $j_2 = 2 \cdot 1 = 2$.
    \item $\text{var}_2 = \text{var}_1 + \frac{8}{12} \cdot 2^2 = 0.6667 + 0.6667 \times 4 = 0.6667 + 2.6667 = 3.3333$.
    \item $\sigma_{\text{ERF}} = \sqrt{3.3333} \approx 1.8257$.
\end{itemize}

Output pixel $i=0$ has center at $0.5$ with receptive span $[-2.5, 3.5]$ (width $7$).

\section{Numerical Considerations and Edge Cases}

\begin{itemize}
    \item \textbf{Asymmetric Padding:} When top/bottom or left/right paddings differ, receptive field calculations must track height and width components $(r_h, r_w)$ and $(\text{start}_h, \text{start}_w)$ independently.
    \item \textbf{Dilated Kernels:} Dilation exponentially expands theoretical RF without adding parameter weights, but can create grid/checkerboard blind spots if consecutive dilations share common factors.
\end{itemize}

\section{References}
\begin{enumerate}
    \item Luo, W., Li, Y., Urtasun, R., & Zemel, R. (2016). Understanding the effective receptive field in deep convolutional neural networks. \emph{NeurIPS}, 4898--4906.
    \item Araujo, A., Norris, W., & Sim, J. (2019). Computing receptive fields of convolutional neural networks. \emph{Distill}, 4(11), e21.
\end{enumerate}

\end{document}
